% GENERATED FILE. Edit canonical lessons, then run npm run generate:books.
% canonical-source-hash: fnv1a64:6a2237a4a13387b6
% canonical-lessons: RU-C322-rytsar, RU-C322-dospekhi, RU-C322-mech, RU-C322-shchit, RU-C322-tron

\chapter{Knights}
\label{ch:ru-knights}
\hlchaptermodality{322}

\begin{chapteropening}
\textbf{By the end of this chapter:} \emph{I can talk about a knight, armour, a sword, a shield and a throne.}

Complete the last lesson of chapter 322: I can talk about a knight, armour, a sword, a shield and a throne.
\end{chapteropening}

\section[\texorpdfstring{\ru{рыцарь} (rýtsar)}{рыцарь (rýtsar)}]{\ru{рыцарь} (rýtsar) — a knight}

\label{lesson:RU-C322-rytsar}



\begin{quote}
Before the new one: say the Russian for a prince, then the Russian for a kingdom.
\end{quote}

\subsection*{You'll want to know: \ru{рыцарь}}
\textbf{\ru{рыцарь}} (\emph{rýtsar}) — \textquotedblleft{}a knight\textquotedblright{}.

\begin{quote}
\ru{Рыцарь} \ru{сел} \ru{на} \ru{коня} \ru{и} \ru{уехал}.
\end{quote}

\begin{quote}
\emph{Rýtsar sel na konyá i uyékhal.}
\end{quote}

\begin{quote}
The knight got on his horse and rode off.
\end{quote}

Genitive \emph{\ru{рыцаря}}.

\begin{cousinweb}
\emph{\ru{рыцарь}} was borrowed through Polish \emph{rycerz} from German \emph{Ritter} \textquotedblleft{}a rider, a knight\textquotedblright{} (EtymDB).
\end{cousinweb}

\begin{grammarlens}[title={What you've built}]
One more word for telling, explaining and describing at B1.
\end{grammarlens}

\subsection*{Your turn}
\begin{itemize}
\raggedright
  \item \emph{Say it:} \emph{rýtsar}
  \item \emph{Say it:} \emph{rýtsar}, once more
  \item \emph{Say it:} a sentence of your own with \emph{rýtsar}
  \item \emph{Recall:} say the Russian for a prince, then the Russian for a kingdom, then say \emph{rýtsar} again
\end{itemize}

\begin{tcolorbox}[breakable,colback=teal!4,colframe=teal!35!black,title={Before you move on}]
What does \ru{рыцарь} mean? (\textquotedblleft{}A knight\textquotedblright{}.) Say it once more.
\end{tcolorbox}
\section[\texorpdfstring{\ru{доспехи} (dospékhi)}{доспехи (dospékhi)}]{\ru{доспехи} (dospékhi) — armour}

\label{lesson:RU-C322-dospekhi}



\begin{quote}
Before the new one: say the Russian for a kingdom, then the Russian for a knight.
\end{quote}

\subsection*{You'll want to know: \ru{доспехи}}
\textbf{\ru{доспехи}} (\emph{dospékhi}) — \textquotedblleft{}armour\textquotedblright{}.

\begin{quote}
\ru{Доспехи} \ru{рыцаря} \ru{были} \ru{очень} \ru{тяжёлыми}.
\end{quote}

\begin{quote}
\emph{Dospékhi rýtsarya býli óchen tyazhyólymi.}
\end{quote}

\begin{quote}
The knight's armour was very heavy.
\end{quote}

Plural; genitive \emph{\ru{доспехов}}.

\begin{cousinweb}
\emph{\ru{доспехи}} is usually traced to \emph{\ru{доспеть}} \textquotedblleft{}to get ready\textquotedblright{}, from the root of \emph{\ru{спех}} \textquotedblleft{}haste\textquotedblright{} and \emph{\ru{успех}} \textquotedblleft{}success\textquotedblright{}: the gear made ready. EtymDB lists no parent.
\end{cousinweb}

\begin{grammarlens}[title={What you've built}]
One more word for telling, explaining and describing at B1.
\end{grammarlens}

\subsection*{Your turn}
\begin{itemize}
\raggedright
  \item \emph{Say it:} \emph{dospékhi}
  \item \emph{Say it:} \emph{dospékhi}, once more
  \item \emph{Say it:} a sentence of your own with \emph{dospékhi}
  \item \emph{Recall:} say the Russian for a kingdom, then the Russian for a knight, then say \emph{dospékhi} again
\end{itemize}

\begin{tcolorbox}[breakable,colback=teal!4,colframe=teal!35!black,title={Before you move on}]
What does \ru{доспехи} mean? (\textquotedblleft{}Armour\textquotedblright{}.) Say it once more.
\end{tcolorbox}
\section[\texorpdfstring{\ru{меч} (mech)}{меч (mech)}]{\ru{меч} (mech) — a sword}

\label{lesson:RU-C322-mech}



\begin{quote}
Before the new one: say the Russian for a knight, then the Russian for armour.
\end{quote}

\subsection*{You'll want to know: \ru{меч}}
\textbf{\ru{меч}} (\emph{mech}) — \textquotedblleft{}a sword\textquotedblright{}.

\begin{quote}
\ru{Рыцарь} \ru{поднял} \ru{меч}.
\end{quote}

\begin{quote}
\emph{Rýtsar pódnyal mech.}
\end{quote}

\begin{quote}
The knight raised his sword.
\end{quote}

Genitive \emph{\ru{меча}}.

\begin{cousinweb}
\emph{\ru{меч}} is inherited from Proto-Slavic \emph{mĭčĭ} (EtymDB), which is usually traced to an early Germanic word for a sword.
\end{cousinweb}

\begin{grammarlens}[title={What you've built}]
One more word for telling, explaining and describing at B1.
\end{grammarlens}

\subsection*{Your turn}
\begin{itemize}
\raggedright
  \item \emph{Say it:} \emph{mech}
  \item \emph{Say it:} \emph{mech}, once more
  \item \emph{Say it:} a sentence of your own with \emph{mech}
  \item \emph{Recall:} say the Russian for a knight, then the Russian for armour, then say \emph{mech} again
\end{itemize}

\begin{tcolorbox}[breakable,colback=teal!4,colframe=teal!35!black,title={Before you move on}]
What does \ru{меч} mean? (\textquotedblleft{}A sword\textquotedblright{}.) Say it once more.
\end{tcolorbox}
\section[\texorpdfstring{\ru{щит} (shchit)}{щит (shchit)}]{\ru{щит} (shchit) — a shield}

\label{lesson:RU-C322-shchit}



\begin{quote}
Before the new one: say the Russian for armour, then the Russian for a sword.
\end{quote}

\subsection*{You'll want to know: \ru{щит}}
\textbf{\ru{щит}} (\emph{shchit}) — \textquotedblleft{}a shield\textquotedblright{}.

\begin{quote}
\ru{На} \ru{щите} \ru{был} \ru{нарисован} \ru{лев}.
\end{quote}

\begin{quote}
\emph{Na shchité byl narisóvan lev.}
\end{quote}

\begin{quote}
A lion was painted on the shield.
\end{quote}

Genitive \emph{\ru{щита}}. Also a board: \emph{\ru{электрический} \ru{щит}} — a fuse box.

\begin{cousinweb}
\emph{\ru{щит}} is inherited through Old East Slavic \emph{ščitŭ} from Proto-Slavic \emph{ščitŭ} \textquotedblleft{}shield\textquotedblright{}, which EtymDB traces to PIE \emph{skey-} \textquotedblleft{}to cut, to split\textquotedblright{}: a split board.
\end{cousinweb}

\begin{grammarlens}[title={What you've built}]
One more word for telling, explaining and describing at B1.
\end{grammarlens}

\subsection*{Your turn}
\begin{itemize}
\raggedright
  \item \emph{Say it:} \emph{shchit}
  \item \emph{Say it:} \emph{shchit}, once more
  \item \emph{Say it:} a sentence of your own with \emph{shchit}
  \item \emph{Recall:} say the Russian for armour, then the Russian for a sword, then say \emph{shchit} again
\end{itemize}

\begin{tcolorbox}[breakable,colback=teal!4,colframe=teal!35!black,title={Before you move on}]
What does \ru{щит} mean? (\textquotedblleft{}A shield\textquotedblright{}.) Say it once more.
\end{tcolorbox}
\section[\texorpdfstring{\ru{трон} (tron)}{трон (tron)}]{\ru{трон} (tron) — a throne}

\label{lesson:RU-C322-tron}



\begin{quote}
Before the new one: say the Russian for a sword, then the Russian for a shield.
\end{quote}

\subsection*{You'll want to know: \ru{трон}}
\textbf{\ru{трон}} (\emph{tron}) — \textquotedblleft{}a throne\textquotedblright{}.

\begin{quote}
\ru{Король} \ru{сидел} \ru{на} \ru{золотом} \ru{троне}.
\end{quote}

\begin{quote}
\emph{Koról sidél na zolotóm tróne.}
\end{quote}

\begin{quote}
The king sat on a golden throne.
\end{quote}

\emph{\ru{Взойти} \ru{на} \ru{трон}} — to come to the throne.

\begin{cousinweb}
\emph{\ru{трон}} came through Latin \emph{thronus} \textquotedblleft{}throne\textquotedblright{} from Greek \emph{thrónos} \textquotedblleft{}a chair, a throne\textquotedblright{} (EtymDB).
\end{cousinweb}

\begin{grammarlens}[title={What you've built}]
One more word for telling, explaining and describing at B1.
\end{grammarlens}

\subsection*{Your turn}
\begin{itemize}
\raggedright
  \item \emph{Say it:} \emph{tron}
  \item \emph{Say it:} \emph{tron}, once more
  \item \emph{Say it:} a sentence of your own with \emph{tron}
  \item \emph{Recall:} say the Russian for a sword, then the Russian for a shield, then say \emph{tron} again
\end{itemize}

\begin{tcolorbox}[breakable,colback=teal!4,colframe=teal!35!black,title={Before you move on}]
What does \ru{трон} mean? (\textquotedblleft{}A throne\textquotedblright{}.) Say it once more.
\end{tcolorbox}
